\section{A finite rational calculator for the inert coefficients}
\label{calc:section}

This section gives a direct computation of $A_m^\infty(s)$, without
enumerating all pruned environment profiles. The two budgets count
different events: $R$ bounds \emph{all} actual core-clock rings, including
null choices of the inert parent, while $N$ bounds real vertex arrivals
after the initial tag and its one leaf. Throughout, $m\ge1$ and $s>0$.

\begin{lemma}[A universal core-ring tail]\label{calc:ring-tail}
Put
\[
 P_R(z)=\prod_{i=1}^R\left(1+\frac zi\right),
 \qquad P_0(z)=1.
\]
The expected discounted generation-$m$ births occurring after the $R$th
actual core-clock ring in the inert process are at most
\begin{equation}\label{calc:tail-formula}
 T_{m,R}(s)=
 \left(\prod_{i=1}^R\frac{i}{i+s}\right)
 \sum_{j=0}^{m-1}[z^j]P_R(z)\,s^{-(m-j)}.
\end{equation}
This is rational for rational $s$. For fixed $m,s$ it tends to zero,
with $T_{m,R}(s)=O_{m,s}(R^{-s}(1+\log R)^{m-1})$.
\end{lemma}

\begin{proof}
Use the common potential genealogy in which every ring of every
potential core creates a potential child, even if its actual counterpart
has a null or nonpromotion event. Actual cores and their clock rings
are subsets of this potential process. The potential core population
is a rate-one Yule process started with one vertex. Write $\tau_R$ for
its $R$th birth time. The $R$th actual core ring occurs no earlier than
$\tau_R$, and every omitted actual generation-$m$ birth is a potential
generation-$m$ birth after $\tau_R$.

At potential size $i$, the holding time has rate $i$, independently of
the uniform choice of its parent. Thus the generation profile at
$\tau_R$ is independent of $\tau_R$, and
\[
 \mathbb E e^{-s\tau_R}=\prod_{i=1}^R\frac{i}{i+s}.
\]
Its expected profile polynomial is $P_R$: adding a vertex at size $i$
multiplies the previous expected polynomial by $1+z/i$.

Conditional on the profile at $\tau_R$, a currently existing vertex
at depth $j<m$ contributes at most, and in the potential process exactly,
$s^{-(m-j)}$ from its new future branches. These contributions do not
double count existing child branches: a future vertex is assigned to
its deepest ancestor already present at $\tau_R$. Summing, and using
independence, proves~\eqref{calc:tail-formula}.
Finally $[z^j]P_R\le H_R^j/j!$, where $H_R=\sum_{i=1}^R1/i$, and the
product in~\eqref{calc:tail-formula} is $O_s(R^{-s})$.
\end{proof}

\subsection{Bounding the histories discarded by a birth budget}

The retained state consists of a finite rooted core, its leaf counts,
and the initial inert stub. Its core depths are the genealogical
depths from the original tag. For such a state $x$, define
\begin{equation}\label{calc:future-reward}
 V_m(x)=\sum_{v:\,d(v)<m}s^{-(m-d(v))},
 \qquad C_m(s)=\max_{0\le r<m}s^{-r}.
\end{equation}
The potential-genealogy comparison shows that $V_m(x)$ bounds all
future generation-$m$ reward from that state, independent of its
leaf counts.

Between core rings, integrate the leaf growth using the exact beta
burst kernel. If a burst is discarded because it exceeds the real-birth
budget, no core birth occurs before its next forced core ring, and its
core graph and depths are unchanged. That ring can contribute at most
one new core to the future reward, or one immediate generation-$m$
birth. Consequently the entire omitted reward from this branch is at
most its discounted omitted-ring mass times
\begin{equation}\label{calc:omitted-burst}
 V_m(x)+C_m(s).
\end{equation}
The omitted-ring mass is computed exactly from the Yule burst tail,
including its factor $k/(k+s)$ for $k$ current cores.

For a retained burst whose subsequent real core event itself exceeds
the birth budget, use the sharper bound
\begin{equation}\label{calc:omitted-core-event}
 \mathbf1_{\{\text{new core has depth }m\}}+V_m(x_{\mathrm{post}}),
\end{equation}
weighted by that event's exact discounted mass. A nonpromotion event
has no immediate reward. Null stub choices use one core ring but no
real birth and are retained accordingly.

\begin{proposition}[Certified scalar intervals]\label{calc:calculator}
Enumerating beta bursts through at most $R$ core rings and at most $N$
real births gives a positive finite rational sum $L_{m,R,N}(s)$ at
rational $s$. Let $E_{m,R,N}(s)$ be the sum of the discarded-history
bounds~\eqref{calc:omitted-burst}--\eqref{calc:omitted-core-event}. Then
\begin{equation}\label{calc:interval}
 L_{m,R,N}(s)\le A_m^\infty(s)
 \le L_{m,R,N}(s)+E_{m,R,N}(s)+T_{m,R}(s).
\end{equation}
The upper endpoint may also be capped at $s^{-m}$. These bounds compute
$A_m^\infty(s)$ to any prescribed rational accuracy: first increase
$R$ until its explicit tail is small, then increase $N$ until the
computed discarded-history bound is small.
\end{proposition}

\begin{proof}
Retained histories have their exact positive discounted weights and
are a subset of the original histories. The first discarded burst or
core event accounts for a missing path by
\eqref{calc:omitted-burst} or~\eqref{calc:omitted-core-event}. The additional
universal tail bounds all births after $R$ actual core rings; any
overlap between the two error terms is harmless. This proves
\eqref{calc:interval}.

For fixed $R$, at most $R+1$ real cores exist through these rings, so
the reward bounds at every first discarded history are bounded by a
constant depending only on $R,m,s$. The first-discard histories are
disjoint. Their total undiscounted probability is at most the
probability that the real population at the $R$th core ring exceeds
$N+2$. That population is finite almost surely: the $R$th core ring
occurs by the $R$th ring of the original tag, and the vertex process
is nonexplosive, with birth rate at most its current population.
Thus $E_{m,R,N}(s)\to0$ as $N\to\infty$ for fixed $R$.

Every step is finite: before each retained ring there are only
finitely many leaf allocations within the birth budget, and null
events consume the ring budget. At rational $s$, each beta entry,
neighbour-choice rate, tail, and reward bound is rational.
\end{proof}

\subsection{Implementation and the complete scalar lower test}

The standard-library program \texttt{inert\_growth\_kernel.py} imports
the exact burst routines from \texttt{core\_clock\_kernel.py}. It keeps
the original tag at root zero and combines root-isomorphic states.
It reports the lower and upper coefficient bounds, the two separate
error terms, and the scalar test
\begin{equation}\label{calc:scalar-test}
 \frac{L_{m,R,N}(s)}{12(s+3)}>1.
\end{equation}
The uniform inert minorization and its growth-limit theorem make this
a sufficient certificate of $s<\mu$. They also make the test complete:
if $s<\mu$, some $m$ has $A_m^\infty(s)>12(s+3)$, and sufficiently
large finite budgets preserve that strict inequality. This is a
semidecision procedure, with no practical search-length estimate near
$\mu$ and no termination claim when $s\ge\mu$.
The optional \texttt{--search} mode implements the enumeration: at stage
$h$, take $R=N=h$ and test generations $1,\ldots,h$. For any fixed
successful generation the retained sums increase to its full coefficient
along these budgets. The option \texttt{--max-stage} imposes a finite
practical stopping point, whose failure gives no conclusion about $s-\mu$.

For example, with just one allowed real birth and $m=1$, the only
rewarded histories are initial null rings followed by a budding ring.
Writing $a=1/[2(s+2)]$, their exact retained sum is
\[
 L_{1,R,1}(s)=\frac{1-a^R}{2s+3}.
\]
The checker verifies this identity separately from its burst and
individual-clock comparisons.

Run, for example,
\begin{quote}\small\ttfamily
python3 scripts/inert\_growth\_kernel.py --check\\
python3 scripts/inert\_growth\_kernel.py --s 1/2 \textbackslash\\
\hspace*{1em}--generation 2 --core-rings 3 --max-births 4
\end{quote}
The checker compares $180$ small burst calculations against an
independent recursion over individual leaf and core clock events,
including null events. It also checks the potential-profile tail
identities, exact one-ring scalar partial sums, and nested certified
intervals. These are finite verification checks, not a new numerical
exponent bound or a replacement for the infinite coupling proofs.
